% Generated by Listas/make_statements.py from pratica-lecture-02.md, item 22. Do not edit.
\textbf{Source.} Tadelis (2013), \emph{Game Theory: An Introduction}, Exercise 5.15 ``Hotelling's Continuous Model'', p.~99 {[}PDF p.~118{]}. Answer in the source: none (the preface, p.~xiii, says about half of the exercises are solved on Princeton UP's website). Reproduced verbatim. \textbf{Reading, not covered in class} (Tadelis §5.2.5, pp.~93--95).



\begin{quote}

\textbf{Hotelling's Continuous Model:} Consider Hotelling's model, in which citizens are a continuum of voters on the interval \(A = [-a, a]\), with uniform distribution \(U(a)\).



\begin{enumerate}

\def\labelenumi{\alph{enumi}.}

\item

  What is the best response of candidate \(i\) if candidate \(j\) is choosing \(a_j > 0\)?

\item

  Show that the unique Nash equilibrium is \(a_1 = a_2 = 0\).

\item

  Show that for a general distribution \(F(\cdot)\) over \([-a, a]\) the unique Nash equilibrium is where each candidate chooses the policy associated with the median voter.

\end{enumerate}

\end{quote}



\textbf{Referred to in the statement.} Hotelling's model (Tadelis §5.2.5, pp.~92--94 {[}PDF p.~111--113{]}) (reproduced from the book):



\begin{quote}

Two politicians, each caring only about being elected; each candidate \(i\) chooses a platform \(a_i\) (the policy of the citizen for whom it is ideal); each citizen votes for the candidate whose platform is closest to his preferences (an indifferent citizen tosses a coin); a candidate wins if a majority votes for him, and candidates prefer winning to a tie and a tie to losing.



\emph{Remark.} In the more common representation of the Hotelling competition model, the citizens are a continuum of voters, given by the interval \(A = [\underline{a}, \overline{a}]\), with distribution \(F(a)\) to determine the distribution of each political preference. We define the median voter as that voter \(a_m\) for which \(\Pr\{a \le a_m\} = \frac12\). For example, if \(F(\cdot)\) is uniform then \(a_m = \frac{\underline{a}+\overline{a}}{2}\).

\end{quote}



\textbf{Optional extras.}

- Tadelis 5.4 (p.~96, splitting pizza) and MWG 8.D.4 (p.~263, splitting 100 dollars, with dominance parts): more continua of equilibria, after item 4.

- Gibbons 1.13 (p.~51, two workers applying to two firms): pure and mixed equilibria in an applied \(2\times2\) game, after item 8.

- Tadelis 6.6 (p.~124, declining industry): a \(3\times3\) game where weak dominance shrinks the support of the unique mixed equilibrium.

- Muñoz-Garcia Ch. 2 Ex. 16 (p.~76, Cournot with efficiency changes) and Ch. 3 Ex. 10 (p.~114, rock-paper-scissors): worked answers in the book.

- Reading, not covered in class: Tadelis 5.7 (p.~96, the \(n\)-player commons), 4.8--4.9 (p.~78, the beauty contest); Gibbons 1.8 (p.~50, voting with two and three candidates); MWG 8.D.5 (pp.~263--264, ice-cream vendors).



\begin{center}\rule{0.5\linewidth}{0.5pt}\end{center}

